\section{Conclusion}
\label{section:conclusion}

We extended the Hyp-RL formulation along two axes---a derivative-informed state representation and cross-dataset meta-training---and evaluated both on a homogeneous learning-curve benchmark (LCBench) and a heterogeneous suite of real-world tabular tasks (Kaggle Playground Series). Our central finding is that the value of the learning-curve derivative features is benchmark-dependent. On LCBench, where every task shares the same network family, search space, and training protocol, the curve-informed \texttt{CrossDataset-LC-DQN} is statistically indistinguishable from its curve-free ablation \texttt{CrossDataset-HyperRL}; the stereotyped, near-monotone curves leave little signal for the derivatives to exploit. On the heterogeneous Kaggle suite, by contrast, the same features produce a consistent cold-start advantage: \texttt{LC-DQN} attains lower average normalized simple regret than \texttt{HyperRL} ($0.0217$ vs.\ $0.0251$) and is the only RL variant to cross the Random Search baseline. Because the two variants differ \emph{only} in the appended derivative statistics, this gap directly attributes the improvement to the learning-curve information.

Two secondary observations qualify the result. First, Bayesian Optimization remains the strongest method asymptotically on both benchmarks; the RL agents' advantage is concentrated in the low-budget cold-start regime, where a transferred prior matters more than an on-the-fly surrogate. Second, cross-dataset transfer does not improve monotonically with the number of meta-training tasks, suggesting that na\"ively pooling dissimilar datasets can inject conflicting reward scales and curve dynamics. Overall, learning-curve derivatives are a cheap, zero-evaluation-cost state augmentation that pays off precisely when curve shape varies across configurations---a property of practical, heterogeneous HPO workloads. This motivates the surrogate-accelerated extension discussed next, which would let the agent exploit this signal across far more datasets without the cost of full training runs.
